% Generated draft: requires mathematical and source review.
\documentclass[11pt]{article}
\usepackage[margin=25mm]{geometry}
\usepackage{fontspec}
\setmainfont{Latin Modern Roman}
\usepackage{amsmath,amssymb,mathtools,unicode-math}
\setmathfont{Latin Modern Math}
\usepackage{enumitem,needspace,xurl,hyperref}
\setlength{\parindent}{0pt}
\setlength{\parskip}{5pt}
\setlength{\emergencystretch}{3em}
\newcommand{\N}{\mathbb{N}}
\newcommand{\Z}{\mathbb{Z}}
\newcommand{\Q}{\mathbb{Q}}
\newcommand{\R}{\mathbb{R}}
\newcommand{\C}{\mathbb{C}}
\newcommand{\F}{\mathbb{F}}
\newcommand{\A}{\mathbb{A}}
\newcommand{\PP}{\mathbb{P}}
\newcommand{\E}{\mathbb{E}}
\newcommand{\eps}{\varepsilon}
\newcommand{\dd}{\,\mathrm{d}}
\begin{document}
\section*{UnsolvedMath ID: 2100401}
\textbf{Source problem code:} catalogue entry
\subsection{Definitions and mathematical statement}
In the smooth strictly convex planar interpretation used by the supplied research summary, let $K\subset\R^2$ be a compact strictly convex table with nonempty interior. For $x\in\R^2\setminus K$, choose one of the two consistently oriented supporting tangents from $x$ to $K$, let $p(x)$ be its tangency point, and define the outer billiard map
\[
T(x)=2p(x)-x.
\]
Choosing the opposite tangent reverses the map. An invariant curve $C$ obeys $T(C)=C$.

The rigidity question asks: if the \emph{entire exterior} $\R^2\setminus K$ is foliated by invariant curves, must $K$ be an ellipse? In the standard integrable-foliation interpretation these are smooth simple closed curves surrounding $K$, forming the leaves of a foliation of its exterior. The source sentence itself does not state their smoothness, closedness, or the regularity of $\partial K$; the research summary explicitly uses smooth invariant curves and a strictly convex planar table. These regularity and foliation conventions require checking in the exact theorem.

The requested conclusion is that there are $c\in\R^2$ and a symmetric positive-definite $2\times2$ matrix $Q$ such that
\[
K=\{x\in\R^2:(x-c)^{\mathsf T}Q(x-c)\leq1\}.
\]
The hypothesis is a global exterior foliation, not merely invariant curves near the table, far from it, or on a subset of positive measure. The supplied summary attributes an affirmative resolution in the smooth strictly convex setting to Bialy--Mironov (2023); no proof or independently checked regularity threshold is supplied here.
\subsection{Short English statement}
The source asks whether a planar outer billiard whose whole exterior is foliated by invariant curves must have an elliptical table. The supplied research summary reports an affirmative theorem for strictly convex tables and smooth invariant curves, with exact regularity hypotheses still to be checked.
\subsection{Sources}
\begin{itemize}
\item Local UnsolvedMath registry, permanent ID 2100401, exact source code catalogue entry, ``Open Problems in Integrable Systems --- Around the Birkhoff conjecture''; status solved, classification SOLVED-IN-LITERATURE.
\item \emph{Open Problems, Questions, and Challenges in Finite-Dimensional Integrable Systems} (2018), Problem 4.1, \url{https://arxiv.org/abs/1804.03737}. The record identifies extraction from public arXiv source TeX, accessed July 29, 2026, and says the source presented this as open.
\item Bialy--Mironov, \emph{Outer billiard rigidity}, arXiv:2306.12494 (2023), cited by the supplied research summary for the affirmative rigidity result. The summary says the paper was announced/available in 2023 and describes it as an outer-billiard analogue of Birkhoff/inner-billiard rigidity results. It supplies no theorem number, precise smoothness assumptions, or refereed publication details.
\end{itemize}
\end{document}
